% Part: computability
% Chapter: computability-theory
% Section: russells-paradox

\documentclass[../../../include/open-logic-section]{subfiles}

\begin{document}

% SEGMENT OLP-0236-S01

\olfileid{cmp}{thy}{rus}
\olsection{ரசலின் முரண்புதிருடன் ஒப்பீடு}

இப்பிரிவிலுள்ள வாதங்களை ரசலின் முரண்புதிருடன் ஒப்பிட்டும்
வேறுபடுத்தியும் பார்ப்பது பயனளிக்கும்:
\begin{enumerate}
\item ரசலின் முரண்புதிர்: $S = \Setabs{x}{x \notin x}$ என்க.
அப்போது $S \in S$ எனில், எனில் மட்டுமே, $X \notin S$; இது ஒரு
முரண்பாடு.

  \emph{முடிவு:} அத்தகைய கணம்~$S$ எதுவும் இல்லை. “அனைத்துக்
  கணங்களின் கணம்” ஒன்று உள்ளது என்று அனுமானிப்பது கணக் கோட்பாட்டின்
  பிற அடிகோள்களுடன் முரணானது.

\item ரசலின் முரண்புதிரின் ஒரு மாற்றம்: அனைத்துச் சார்புகளின்
  கணத்திலிருந்து $\{ 0, 1 \}$ க்குச் செல்லும் “சார்பு” $F$
  பின்வருமாறு வரையறுக்கப்படட்டும்:
  \[
  F(f) =
  \begin{cases}
    1 & \text{$f$ ஆனது $f$ இன் சார்பகத்தில் இருந்து, $f(f) = 0$ எனில்} \\
    0 & \text{இல்லையெனில்}
  \end{cases}
  \]
  $F(F) = 0$ எனில், எனில் மட்டுமே, $F(F) = 1$ என்று ஒத்த
  வாதம் காட்டுகிறது; இது ஒரு முரண்பாடு.

  \emph{முடிவு:} $F$ ஒரு சார்பு அல்ல. “அனைத்துச் சார்புகளின் கணம்”
  ஒரு சார்பின் சார்பகமாக இருக்க முடியாத அளவு பெரியது.

\item மூலைவிட்டமாக்கல் வாதம்: $f_0$, $f_1$, \dots என்பது பகுதியாகக்
  கணிக்கத்தக்க சார்புகளின் எண்ணிடல் ஆகட்டும்; மேலும்
  $G \colon \Nat \to \{ 0, 1 \}$ பின்வருமாறு வரையறுக்கப்படட்டும்:
  \[
  G(x) =
  \begin{cases}
    1 & \text{$f_x(x)\downarrow = 0$ எனில்} \\
    0 & \text{இல்லையெனில்}
  \end{cases}
  \]
  $G$ கணிக்கத்தக்கது எனில், அது ஏதேனும் $k$ க்கான சார்பு $f_k$.
  ஆனால் அப்போது $G(k) = 1$ எனில், எனில் மட்டுமே, $G(k) = 0$;
  இது ஒரு முரண்பாடு.

  \emph{முடிவு:} $G$ கணிக்கத்தக்கதல்ல. கணக் கோட்பாட்டின்
  அடிகோள்களின்படி $G$ இன்னும் ஒரு சார்பே என்பதைக் கவனிக்க;
  இங்கே முரண்புதிர் எதுவும் இல்லை, ஒரு தெளிவுபடுத்தல் மட்டுமே.
\end{enumerate}

பகுதிச் சார்புகள், கணிக்கத்தக்க சார்புகள், பகுதியாகக் கணிக்கத்தக்க
சார்புகள் முதலியவற்றைப் பற்றிய பேச்சு குழப்பமளிக்கலாம்.
$\Nat$ இலிருந்து $\Nat$ க்குச் செல்லும் அனைத்துப் பகுதிச்
சார்புகளின் கணம் ஒரு பெரிய பொருட்திரட்டு. அவற்றுள் சில முழுச்
சார்புகள், சில கணிக்கத்தக்கவை, சில முழுச் சார்புகளாகவும்
கணிக்கத்தக்கவையாகவும் உள்ளன, சில இரண்டுமல்ல. “சார்பு” என்று
மட்டும் கூறும்போது, இயல்பாக முழுச் சார்பையே குறிக்கிறோம் என்பதை
மனத்தில் கொள்க. எனவே நம்மிடம் பின்வருவன உள்ளன:
\begin{enumerate}
\item கணிக்கத்தக்க சார்புகள்
\item முழுச் சார்புகள் அல்லாத பகுதியாகக் கணிக்கத்தக்க சார்புகள்
\item கணிக்கத்தக்கவை அல்லாத சார்புகள்
\item முழுச் சார்புகளாகவும் கணிக்கத்தக்கவையாகவும் இல்லாத பகுதிச்
சார்புகள்
\end{enumerate}
இதை ஒழுங்குபடுத்த, $\Nat$ இலிருந்து $\Nat$ க்குச் செல்லும் அனைத்துப்
பகுதிச் சார்புகளையும் குறிக்கும் ஒரு பெரிய சதுரத்தை வரைந்து,
முறையே முழுச் சார்புகளுக்கும் கணிக்கத்தக்க பகுதிச் சார்புகளுக்கும்
ஒத்த இரு ஒன்றன்மேலொன்று படியும் பகுதிகளைக் குறித்துக் காட்டுவது
உதவலாம். படத்தில் இவ்வாறு கிடைக்கும் ஒவ்வொரு பகுதியிலும் ஓர்
எடுத்துக்காட்டை விவரிக்க முடியுமா என்று பார்ப்பது ஒரு நல்ல பயிற்சி.

\end{document}
